% 제4장: 실험 결과 및 분석 / 4.2 초안
% section4_1.tex 다음에 삽입한다. 표는 4.4에서도 함께 참조한다.
% 근거: results/koblex-reproduction-v1/analysis.json,
%       results/koblex-reproduction-v1/chapter4-prep-v1/A1/numeric-ledger.json,
%       results/koblex-reproduction-v1/chapter4-prep-v1/A4/stage-summary.jsonl,
%       results/koblex-reproduction-v1/chapter4-prep-v1/A4/pool-feasibility-summary.json

\subsection{재정렬 모델 교체에 따른 근거 선택 성능}
\label{sec:ch4_evidence_results}

표~\ref{tab:ch4_main_results}는 전체 226문항의 최종 근거 선택 성능과
답변 품질을 제시한다. 근거 선택 성능은 공식 정답 근거(gold)와
LLM이 최종 선택한 근거를 문자열 집합으로 비교하고, 문항별 점수를
비가중 평균하여 산출하였다. EM은 두 집합의 완전 일치 여부를,
모든 gold 포함률은 추가 근거의 유무와 관계없이 정답 근거를 모두
포함한 문항의 비율을 나타낸다.

\begin{table}[htbp]
    \centering
    \caption{수정 CE와 Jev의 전체 성능 비교 ($n=226$)}
    \label{tab:ch4_main_results}
    \small
    \setlength{\tabcolsep}{3pt}
    \begin{tabular}{@{}lrrrr@{}}
        \hline
        지표 & 수정 CE & Jev & $\Delta$ & 차이의 95\% CI \\
        \hline
        \multicolumn{5}{@{}l}{\textbf{근거 선택}} \\
        F1              & 58.57 & 61.75 & $+3.19$ & $[+0.75,\,+5.68]$ \\
        Recall          & 57.74 & 60.69 & $+2.95$ & $[+0.39,\,+5.51]$ \\
        EM              & 27.88 & 29.65 & $+1.77$ & $[-1.77,\,+5.19]$ \\
        모든 gold 포함률 & 34.07 & 34.96 & $+0.88$ & $[-2.90,\,+4.53]$ \\
        \hline
        \multicolumn{5}{@{}l}{\textbf{답변 품질}} \\
        Token F1        & 46.20 & 47.31 & $+1.12$ & $[+0.02,\,+2.27]$ \\
        LF-Eval         & 67.57 & 70.15 & $+2.58$ & $[+0.37,\,+4.79]$ \\
        \hline
    \end{tabular}
    \par\smallskip
    \begin{minipage}{\linewidth}
        \footnotesize
        값은 0--1 원점수의 100배이며, EM과 모든 gold 포함률은 백분율이다.
        $\Delta$는 Jev$-$수정 CE로, EM과 모든 gold 포함률의 차이는 \%p,
        나머지 지표의 차이는 점 단위이다. 차이와 CI는 반올림 전 값으로
        계산하였다. CI는 151개 의존 그룹의 대응 bootstrap 10,000회로
        산출한 탐색적 percentile 구간이며, 다중 비교를 보정하지 않았다.
    \end{minipage}
\end{table}

Jev의 근거 F1은 61.75점으로 수정 CE의 58.57점보다 3.19점 높았으며,
차이의 95\% 신뢰구간은 $[+0.75,\,+5.68]$점이었다.
Recall도 57.74점에서 60.69점으로 2.95점 증가하였고,
차이의 신뢰구간은 $[+0.39,\,+5.51]$점이었다.
두 지표의 탐색적 신뢰구간이 모두 0보다 높다는 결과는
Jev 교체 후 최종 근거 집합에서 정답 근거를 부분적으로 회수하는
정도가 평균적으로 개선되는 방향을 보여준다.

반면, 근거 집합 전체의 완결성에서는 개선의 불확실성이 남았다.
EM의 증가폭은 1.77\%p였으며 신뢰구간이 0을 포함하였다.
모든 gold를 포함한 문항은 수정 CE 77개에서 Jev 79개로 증가하였으나,
비율 차이는 0.88\%p이고 신뢰구간은 $[-2.90,\,+4.53]$\%p였다.
따라서 평균 F1과 Recall의 증가만으로 여러 정답 근거를 모두 확보하는
성능까지 개선되었다고 확정하기는 어렵다.

이 차이를 단계별로 살펴보면, 공통 BM25 후보 합집합의 문항 평균
gold Recall은 77.21\%, 모든 gold 포함률은 61.06\%(138/226)였다.
예상 법령별 상위 10개를 합친 shortlist의 평균 gold Recall은
수정 CE 69.47\%, Jev 74.63\%로, Jev가 더 많은 정답 근거를
후속 선택 단계에 남겼다. 다만 각 후보 목록에서 하나만 선택하는
제약을 고려하면, shortlist에서 모든 gold를 동시에 선택할 수 있는
문항은 수정 CE 104개, Jev 120개였다. 이 값은 구조적 가능성의 상한이며,
실제 모든 gold 선택 문항 수인 77개와 79개에는 후속 LLM 선택도
영향을 미친다. 최종 성능 차이는 재정렬과 공통 LLM 선택이 연결된
파이프라인의 결과로 해석한다.

문자열 매핑의 제약도 함께 고려해야 한다. 고유 gold 386개 중 27개는
corpus에 정확히 매핑되지 않아 29문항에 영향을 주었다.
완전 매핑된 197문항의 F1은 수정 CE 64.12점, Jev 67.21점이었고,
매핑 불일치의 영향을 받은 29문항에서는 각각 20.80점과 24.71점이었다.
두 집단 모두 Jev의 점추정이 높았으나, 전체 결과는 해당 29문항을
유지한 분모로 보고하였다. 또한 exact gold 이외의 선택 근거가 갖는
법률적 타당성은 이 문자열 일치 점수만으로 판단하지 않는다.

한편 공개 코드 CE의 F1은 48.72점으로, 인덱스 문제를 수정한 CE와
9.85점의 차이를 보였다. 이는 목록 참조 인덱스 수정의 영향으로
Jev 교체에 따른 3.19점 차이와 구분한다. 공개 코드 스냅샷에서
확인한 이 문제를 원 논문의 실제 실행에도 존재했던 것으로
단정하지 않는다.
